%% ============================================================
%%  CSCI 4922/5922 — Lab Assignment 3 Report Template
%%  Neural Networks and Deep Learning, Spring 2026
%% ============================================================

\documentclass[11pt]{article}

% ── Packages ────────────────────────────────────────────────
\usepackage[margin=1in]{geometry}
\usepackage{amsmath, amssymb}
\usepackage{booktabs}        % professional tables
\usepackage{graphicx}        % figures
\usepackage{hyperref}        % clickable references
\usepackage{caption}
\usepackage{subcaption}      % subfigures
\usepackage{enumitem}        % tighter lists
\usepackage{titlesec}        % section formatting
\usepackage{parskip}         % space between paragraphs
\usepackage{listings}        % code in appendix
\usepackage{xcolor}
\usepackage{microtype}       % better typography

% ── Code listing style ──────────────────────────────────────
\lstset{
  language=Python,
  basicstyle=\ttfamily\footnotesize,
  keywordstyle=\color{blue},
  commentstyle=\color{gray},
  stringstyle=\color{orange!80!black},
  breaklines=true,
  frame=single,
  numbers=left,
  numberstyle=\tiny\color{gray},
  xleftmargin=1.5em,
}

% ── Convenience command for placeholders ────────────────────
\newcommand{\todo}[1]{\textcolor{red}{\textbf{[TODO: #1]}}}

% ── Title block ─────────────────────────────────────────────
\title{
  \textbf{Lab Assignment 3: Visual Question Answering Tasks}\\
  \large CSCI \todo{4922 or 5922} — Neural Networks and Deep Learning
}
\author{
  \todo{Your Full Name} \\
}
\date{}

% ════════════════════════════════════════════════════════════
\begin{document}
\maketitle

% ════════════════════════════════════════════════════════════
\section{Custom Multi-Modal Architectures (Challenges 1 \& 2)}
% Rubric: 17.5 points total
% ════════════════════════════════════════════════════════════

% ── 1.1 Methods ─────────────────────────────────────────────
\subsection{Methods}

\subsubsection{Dataset Preprocessing}

\todo{
  Describe your preprocessing pipeline in full detail. Cover at minimum:
  (1) image preprocessing (resize dimensions, normalization statistics),
  (2) text preprocessing / tokenization scheme,
  (3) how you handle the binary label for Task 1 and the 10 human answers
      for Task 2 (e.g., majority vote, all-answers supervision, etc.),
  (4) how you sampled up to 10,000 training examples from the full 20,523.
}

\subsubsection{Experiment Design — Task 1: Answerability Classification}

\todo{
  Clearly state: (1) the hypothesis being tested, (2) the experimental
  variable (what changes between conditions), (3) everything held constant,
  and (4) how you will evaluate the outcome (metric, dataset split).
}

\subsubsection{Experiment Design — Task 2: Answer Generation}
% Same rubric criteria as above for the generation task.

\todo{
  Same structure as 1.1.2 but for the text generation task. State the
  hypothesis, variable, controls, and evaluation metric.
}

\subsubsection{Final Architecture Descriptions}

\paragraph{Challenge 1 — Classification Architecture.}
\todo{
  Describe the layers, attention operations, input/output dimensions,
  activation functions, and any other notable design choices of your
  final Challenge 1 model.
}

\paragraph{Challenge 2 — Answer Generation Architecture.}
\todo{
  Same as above for your Challenge 2 model.
}

% ── 1.2 Results ─────────────────────────────────────────────
\subsection{Results}
% Rubric: "Well-formatted results are provided for each experiment" (5 pts)

\subsubsection{Experiment Results — Task 1: Answerability Classification}

\todo{
  Report validation accuracy (or other chosen metric) for each experimental
  condition. Use tables, figures or whatever you believe will be the best way to communicate your results.
}

\todo{This table is just an example, please add your own metrics to report and add them to the table.}
\begin{table}[h]
  \centering
  \caption{Task 1 experiment results: \todo{experiment name, e.g., attention type comparison}.}
  \label{tab:exp1_cls}
  \begin{tabular}{lcc}
    \toprule
    \textbf{Condition} & \textbf{Val Accuracy (\%)} & \textbf{Notes} \\
    \midrule
    \todo{Condition A}  & \todo{XX.X}  & \\
    \todo{Condition B}  & \todo{XX.X}  & \\
    \todo{Condition C}  & \todo{XX.X}  & \\
    \bottomrule
  \end{tabular}
\end{table}

\subsubsection{Experiment Results — Task 2: Answer Generation}

\todo{
  Report validation accuracy (human-answer metric) for each experimental
  condition. Use tables, figures or whatever you believe will be the best way to communicate your results.
}

\begin{table}[h]
  \centering
  \caption{Task 2 experiment results: \todo{This is just a placeholder table. Please adjust the results to whatever you feel is best to articulate what you found.}.}
  \label{tab:exp1_gen}
  \begin{tabular}{lcc}
    \toprule
    \textbf{Condition} & \textbf{Val Answer Accuracy} & \textbf{Notes} \\
    \midrule
    \todo{Condition A}  & \todo{X.XXX}  & \\
    \todo{Condition B}  & \todo{X.XXX}  & \\
    \todo{Condition C}  & \todo{X.XXX}  & \\
    \bottomrule
  \end{tabular}
\end{table}

% ── 1.3 Analysis ────────────────────────────────────────────
\subsection{Analysis}

\subsubsection{Analysis — Task 1 Experiment}
\todo{
  Interpret the results from Table~\ref{tab:exp1_cls}. Which condition
  performed best and why? Were there surprising outcomes? What would you
  do differently with more compute?
}

\subsubsection{Analysis — Task 2 Experiment}
\todo{
  Same depth of analysis for Table~\ref{tab:exp1_gen}.
}


% ════════════════════════════════════════════════════════════
\section{Foundation Model Features — CLIP (Challenges 3 \& 4)}
% Rubric: 12.5 points total
% ════════════════════════════════════════════════════════════

\subsection{Methods}

\iffalse
\subsubsection{Feature Loading and Preprocessing}

\todo{
  Describe how you loaded the CLIP features (512-d image and text vectors),
  any normalization or projection applied before feeding into your model,
  and how you enforced the ≤10,000 training sample limit.
}
\fi

\subsubsection{Experiment Design — Task 1: Answerability Classification (CLIP)}

\todo{
  State hypothesis, experimental variable, controls, and evaluation
  metric for your CLIP-based classification experiment.
}

\subsubsection{Experiment Design — Task 2: Answer Generation (CLIP)}

\todo{
  Same structure for the CLIP-based answer generation experiment.
}

\subsubsection{Final Architecture Descriptions}

\paragraph{Challenge 3 — CLIP Classification Architecture.}
\todo{Describe the architecture built on top of CLIP features.}

\paragraph{Challenge 4 — CLIP Answer Generation Architecture.}
\todo{Describe the architecture built on top of CLIP features.}

\subsection{Results}

\subsubsection{Experiment Results — Task 1 (CLIP)}

\begin{table}[h]
  \centering
  \caption{Task 1 CLIP experiment results: \todo{This is just a placeholder table. Please adjust the results to whatever you feel is best to articulate what you found.}.}
  \label{tab:exp2_cls}
  \begin{tabular}{lcc}
    \toprule
    \textbf{Condition} & \textbf{Val Accuracy (\%)} & \textbf{Notes} \\
    \midrule
    \todo{Condition A}  & \todo{XX.X}  & \\
    \todo{Condition B}  & \todo{XX.X}  & \\
    \bottomrule
  \end{tabular}
\end{table}

\subsubsection{Experiment Results — Task 2 (CLIP)}

\begin{table}[h]
  \centering
  \caption{Task 2 CLIP experiment results: \todo{This is just a placeholder table. Please adjust the results to whatever you feel is best to articulate what you found.}.}
  \label{tab:exp2_gen}
  \begin{tabular}{lcc}
    \toprule
    \textbf{Condition} & \textbf{Val Answer Accuracy} & \textbf{Notes} \\
    \midrule
    \todo{Condition A}  & \todo{X.XXX}  & \\
    \todo{Condition B}  & \todo{X.XXX}  & \\
    \bottomrule
  \end{tabular}
\end{table}

\subsection{Analysis}

\subsubsection{Analysis — Task 1 Experiment (CLIP)}
\todo{Interpret results from Table~\ref{tab:exp2_cls}. Why did certain conditions outperform others?}

\subsubsection{Analysis — Task 2 Experiment (CLIP)}
\todo{Interpret results from Table~\ref{tab:exp2_gen}.}

\subsubsection{End-to-End vs.\ CLIP Feature Models}

\todo{
  Compare the best end-to-end models (Section 1) against the best CLIP
  feature models (Section 2) on both tasks. Which approach performs
  better? Provide a hypothesis for why. Discuss trade-offs in terms of
  model complexity, training time, and final performance.
}


% ════════════════════════════════════════════════════════════
\section{Data Efficiency (CSCI 5922 Only)}
% Rubric: 10 points
% 4922 students: leave this section blank — it will not be graded.
% ════════════════════════════════════════════════════════════

\subsection{Hypothesis}

\todo{
  Before retraining the models, provide a hypothesis on the influence of the extra
training samples. Will the from-scratch model benefit more from the additional
data than the CLIP features model? Why or why not?
}

\subsection{Methods}

\todo{
  List all fixed hyperparameters used for both models (learning rate,
  batch size, number of epochs, optimizer, scheduler, etc.). Explain how
  you constructed the 75\% training split (e.g., random seed, any
  stratification). Note that hyperparameters are frozen from the
  competition models.
}

\subsection{Results}

\todo{This is just a placeholder table. Please adjust the results section to whatever you feel is best to articulate your findings.}

\begin{table}[h]
  \centering
  \caption{Validation accuracy across training data fractions for both binary
           classification architectures.}
  \label{tab:data_eff}
  \begin{tabular}{llccc}
    \toprule
    \textbf{Model} & \textbf{Task} & \textbf{50\% (baseline)} & \textbf{75\%} & \textbf{100\%} \\
    \midrule
    End-to-End  & Classification & \todo{XX.X} & \todo{XX.X} & \todo{XX.X} \\
    CLIP-based  & Classification & \todo{XX.X} & \todo{XX.X} & \todo{XX.X} \\
    \bottomrule
  \end{tabular}
\end{table}

\subsection{Analysis}

\todo{
  Analyze the trends in Table~\ref{tab:data_eff}. Does performance
  increase monotonically with more data? Is the improvement larger for
  one model? Why? Then explicitly revisit your hypothesis from
  Section~\ref{sec:hypothesis}: was it supported, partially supported,
  or refuted? Discuss what this reveals about data efficiency in the
  context of pre-trained vs.\ from-scratch models.
}

\label{sec:hypothesis}


% ════════════════════════════════════════════════════════════
\clearpage
\appendix
\section*{Appendix: Code}
% ════════════════════════════════════════════════════════════
\todo{Please add the code used for this lab here. An example is provided below.}

\begin{verbatim}
// This is a simple code snippet
int main() {
    printf("Hello, world!");
    return 0;
}
\end{verbatim}

\end{document}